\chapter{Accenni al funzionamento di un calcolatore}
Un calcolatore è generalmente composto da:
\begin{itemize}
  \item la $CPU$, esegue le istruzioni, tiene in memoria ($Instruction Register$)
    l'istruzione corrente e la posizione dell'istruzione nella RAM ($P.C.$);
  \item la $RAM$, tiene in memoria le istruzioni da eseguire, certe sezioni
    sono riservate per specifiche funzioni ($stack$, $SP$, $BP$);
  \item dei moduli, come un disco fisso o altro;
  \item la $NIC$, che permette di inviare e ricevere pacchetti in rete.
\end{itemize}

\section{Networking}
Vengono effettuate delle syscall, per eseguire operazioni necessario all'invio
di dati.
Per esempio la syscall send, copia dalla ram una certa sezione di memoria nel
NIC, che verrà inviata in rete.

\subsection{TCP/UDP}
\begin{figure}[t]
\includegraphics[width=.8\textwidth]{gfx/chadudp.jpg}
\centering
\end{figure}

Il protocollo TCP garantisce l'invio del messaggio, dato che attende una
risposta da parte del ricevente.
Il protocollo UDP invia datagrammi, ma non si è certi se la ricezione è
avvenuta.

\subsection{Device driver}
E' un modulo che gestisce le risore fisiche dei vari moduli, tra cui NIC,
gestiscono le interruzioni del dispositivo.

\section{Trasmissione di dati}
\subsection{Invio di dati}
Il DMA presente nel NIC è il componente che accede alla memoria.

La maggior parte dei calcolatori include l'unità di gestione della memoria MMU,
che permette al processore di gestire indirizzi virtuali per ottenere una
maggiore efficienza, ma questo dispositivo è esclusivo del processore, motivo
per cui il DMA del NIC utilizzerà indirizzi fisici.

Delle porzioni di memoria della RAM saranno destinate ad essere utilizzate per
i \emph{buffer}, ovvero delle aree dove NIC può immagazzinare i dati da inviare
o da ricevere.

Al momento dell'invio di un datagramma, esso verà salvato in un buffer, quindi
inviato in rete appena possibile, quindi cancellato dalla memoria.

\subsection{Ricezione di dati}
Per la ricezione di datagrammi dei buffer devono essere sempre disponibili,
cosicchè NIC possa utilizzarli una volta ricevuto il pacchetto.

Il datagramma ricevuto sarà inserito dal DMA in un buffer, un interrupt nel
frattempo creeà nuovi buffer per possibili nuovi messaggi, quindi il sistema
legge il buffer copiandolo nel buffer dell'applicazione, direzionando l'output
verso l'utente corretto, quindi la memoria viene liberata.

Il comportamento standard di un applicazione che implementa networking 
implementa un meccanismo bloccante per cui l'applicazione rimane in attesa
fintanto che il buffer non è stato scritto al suo interno.

\subsection{Trasmissione di dati su bus ring based}
La problematica principale del DMA di tipo \emph{bus mastering} è la
necessità di dover programmare in anticipo i buffer allocati, limitandone
la capacità.

Fintanto che il DMA è impegnato in un operazione perchè in attesa del
processore verrano persi i dati nel frattempo ricevuti dalla rete.

Per risolvere questo problema si implementa una struttura dati più complicata,
integrando una piccola CPU all'interno della NIC (ring based).

\subsubsection{Ring}
Il ring è composto da un anello contenente più di un buffer, e due puntatori,
un puntatore di inizio ed uno di fine gestiti da due processori diversi;
la CPU si occupa di aggiungere elementi all'interno del ring quindi gestirà
l'indice di fine, il puntatore al primo elemento viene gestito dal processore
all'interno della NIC.
SI ottiene una coda di buffer che si possono aggiungere o togliere a seconda
se si sta gestendo un invio o ricezione di dati.

Quindi la NIC legge i buffer in memoria e usa un flag per impostare se il
buffer è libero oppure se deve essere processato.

L'unico tipo di comunicazione tra processore e DMA è la trasmissione della 
quantità di buffer che dovranno essere utilizzati al DMA.

\subsection{Commutazione}
Si rende necessario, per una questione di efficienza, un circuito fisico $R$
attraverso cui trasmettere i pacchetti.

I pacchetti sono inviati come datagrammi, e nella struttura includono il
destinatario e le informazioni trasmesse.

\subsection{Dataflow}
L' \emph{host} compila ed invia il buffer in locale al buffer di ricezione di
$R1$(router), il cui compito è di leggere il datagramma, ricavarne il
destinatario, quindi di inviarlo: l'hos può quindi cancellare il buffer.

La parte ricevente ($R2$) copia il datagramma nel suo buffer di ricezione,
quindi $R1$ può rimuovere dal proprio buffer il datagramma.
Quindi $R2$ invia il datagramma all'effettivo destinatario con un altro
\emph{store and forward}.

Ogni passaggio di trasmissione del datagramma viene chiamato \emph{hop}.

\subsection{Routing}
Identifichiamo per il processo di routing l' \emph{host}, gli endpoint della
trasmissione di pacchetti e i \emph{router}, calcolatori.
Più router possono essere collegati tra loro.

\section{Internet protocol stack}
\begin{enumerate}
  \item Physical
  \item Datalink
  \item Network
  \item Transport
  \item Application
\end{enumerate}

\subsection{Protocollo fisico}
L'unico protocollo fisico sopravissuto ed in uso è il protocollo ethernet.

Ogni richiesta sarà composta da: header, payload e trailer.

Il protocollo definisce che ad ogni host venga assegnato un indirizzo
\emph{MAC}, composto da $6byte$.

L'header conterrà come dati gli indirizzi IP e le porte dei nodi comunicanti.

Il trailer è necessario per il controllo di integrità nella rete ethernet, l'
algoritmo utilizzato è \textbf{CRC32}.

\emph{Il datalink utilizza le stesse modilità di funzionamento del protocollo
fisico}.

\subsection{Network protocol}
Il protocollo network utilizza come gestore degli indirizzi i protocolli
\emph{IPv4 e IPv6}.
Ipv4 utilizza $4byte$ per il funzionamento, significa che può gestire fino a 
circa $4$ miliardi di indirizzi, mentre ipv6 ne può utilizzare molti di più.

I primi sviluppatori dei protocolli non avevano la possibilità
di interagire direttamente con i sistemi operativi, motivo per cui si è dovuto
introdurre il concetto di \textbf{porta di comunicazione}, un numero intero
in $16bit$, con cui si identifica univocamente un numero di una porta ad una
tipologia di applicazione.

Le porte da $0$ a $1023$ sono state dedicate ai protocolli internet ed alle
applicazioni più comuni, un altra serie di porte sono sempre utilizzate per
il networking  mentre le restanti sono le così dette porte effimere che vengono
utilizzate dalle applicazioni in locale.

Quindi se la porta è la parte che un calcolatore espone ad internet, il
\emph{socket} è il corrispettivo componente nel nodo in locale, così da
associare alla richiesta una macchina ed applicazione in particolare.


Sono usati i protocolli \emph{UDP} e \emph{text}
Il protocollo \emph{UDP} invia datagrammi mentre il protocollo \emph{TCP}
permette di implementare una modalità di comunicazione più avanzata, di tipo
\emph{stream}, che permette di interagire con l'interfaccia di comunicazione
come se fossero file.

Gli indirizzi \textbf{IP} sono utilizzati per essere associati agli indirizzi
\textbf{MAC} nel network.

IL \emph{TTL} (Time To Live) definisce da quanto tempo un pacchetto è nel
network.

\emph{RFC} (Remote Function Call) cerca prima di chiamare le versioni dei
protocolli più recenti prima di doverne eseguire uno più vecchio.

\subsubsection{Indirizzi}
Ad ogni macchina corrisponde un IP ($IPv4$ o $IPv6$)
\begin{figure}[H]
    \centering
    \incfig{indirizzi}
    \caption{indirizzi}
    \label{fig:indirizzi}
\end{figure}

Un indirizzo $IPv4$ in default è composto da:
\begin{itemize}
  \item $29bit$ dedicati al network;
  \item $3bit$ dedicati all'host.
\end{itemize}
Per ovviare a problemi di carenza di indirizzi IP, si introduce il concetto
di \emph{NETMASK}, una parola composta da una serie di $1$ ed una di $0$, gli
$1$ rappresentano quanti bit sono dedicati al network, gli $0$ per l'host.

\paragraph{Indirizzi privati}
Sono definiti nel protocollo \emph{RFC 1918} dichiara che un indirizzo
\textbf{IPv4} deve essere composto da $4$ numeri decimali seguiti da punto,
inoltre dichiara che gli indirizzi che cominciano con $10.$ appartengono alla
\emph{classe A}, questi indirizzi sono i cosiddetti indirizzi privati,
utilizzabili solo in \textbf{LAN}.

Gli indirizzi di classe \textbf{B} cominciano da $192.168.0.0$ per arrivare
fino $192.168.255.255$, sono sempre indirizzi privati.


Si può configurare il router per usare il \textbf{NAT} (Network Access
translator), con lo scopo di assegnare due indirizzi IP allo stesso router,
con lo scopo di esporne uno alla rete locale e l'altro a quella internet.

Con il NAT è possibile convertire indirizzi pubblici in indirizzi privati,
direzionando i messaggio alla macchina corretta sostituendo il proprio
indirizzo alla richiesta assegnando una porta libera.

\subsection{Trasporto}
Vengono utilizzato i protocolli \emph{UDP} e \emph{TCP} e la tecnica della
\emph{3 way handshake}.

\begin{figure}[ht]
    \centering
    \incfig{3some}
    \caption{Three way handshake}
    \label{fig:3some}
\end{figure}

Il trasporto comincia con l'invio di un messaggio \textbf{SYN} (accompagnato
da \textbf{SEQ}), all'arrivo
il server risponde con un messaggio che viene chiamato \textbf{SYN-ACK}, la
\textbf{SEQ} e la \textbf{ACK};
quindi l'host invia un uteriore messaggio chiamato \textbf{ACK}, oltre che
i dati.

L'header TCP contiene una serie di informazioni tra le quali i \textbf{FLAG}
SYN e ACK oltre che la porta sorgente (\textbf{SPORT}) e di destinazione
(\textbf{DPORT}).


Vengono utilizzate $2$ parole, la \textbf{SEQ} e la \textbf{ACK}, permettono
la di realizzare l'astrazione dello stream di byte; questi numeri per motivi
di sicurezza non partono da $0$, ma da un valore casuale; lo scopo della
prima squenza di passaggi è di scambiare queste informazioni.

Inoltre \textbf{SEQ} segnala anche l'inizio dello stream di byte da leggere.

Se il messaggio viene suddiviso in tanti datagrammi cumulativi, viene inviato un unico \textbf{ACK cumulativo}.

Il \textbf{timeout} è il tempo che il processo impiega ad inviare e ricevere i
messaggi, dipende dalla lunghezza dei messaggi e dalla velocità dei
calcolatori.

La \textbf{media mobile} dà più peso agli ultimi valori, si rende necessaria
per calcolare il tempo di timeout.
\[
  \text{Estimated }=(1-\alpha)\times\text{ Estimated}+\alpha\text{ SampleRTT}
\]

Inoltre si usa una stime della variabilità:
\[
  \text{Dev}=(1-\beta)\times\text{ Dev}+\beta\times\mid\text{SampleRTT}-
  \text{Estimated}\mid
\]

I parametri $\alpha,\beta$ hanno un valore consigliato: $\alpha=\dfrac{1}{8}$,
$\beta=\dfrac{1}{4}$.

Quindi inizialmente viene assegnato un valore fisso viene inizialmente
assegnato ad \emph{Estimated} e \emph{Dev}, quindi il primo messaggio viene
inviato, quindi nel primo scambuio, in cui il client riceve il \emph{SYN-ACK},
si ottiene il primo \emph{SampleRTT}.

Quindi si possono comparare i valori \emph{Estimated} e \emph{Dev},
(Deviazione) ogni scambio di \emph{ACK e SIN-ACK}.

Il \textbf{timeout sarà}:
\[
  \text{Time-out}=\text{Estimated}+4\times\text{ Dev}
\]

\subsubsection{Flow control}
Il motivo principale della perdita di messaggi è la mancanza di
sincronizzazione, un tempo era molto probabile ricevere messaggi corrotti,
ora la perdita di messaggi è legata alla mancanza di buffer liberi di
ricezione.

Il \textbf{mittente} avrà un buffer di ricezione (TX buffer), così come
il destinatario (RX buffer), nei quali i datagramm vengono ricevuti.

Per garantire che vi siano sempre buffer di ricezione liberi sarà necessario
comunicare all'interlocutore il numero di buffer allocati.

All'interno dell'header \textbf{TCP} vi sarà un campo \textbf{receive window}
che contiene la grandezza dei buffer di ricezione del destinatario.

I messaggi saranno inviati o alla ricezione di un \textbf{ACK} o alla scadenza
del \textbf{TIMEOUT}.

In generale, \emph{UDP} utilizza un approccio ottimistico per cui in una
comunicazione fluida è più efficiente rispetto a \emph{TCP} che si accerta
che ogni messaggio venga ricevuto correttamente.

\subsubsection{Congestion control}
Il \textbf{TCP} implementa anche un controllo di congestione oltre al \emph{
flow control}, è necessario per ridurre la perdita di messaggi sulla reta,
ma non da parte del destinatario, ma dai \emph{router} che implementano lo
\emph{store and forward}.

Quindi è il \emph{router} che rileva la congestione di dati (mancanza di
buffer) e rallenta la ricezione di dati.

Il congestion control avviene in tre fasi:
\begin{enumerate}
  \item \textbf{slow start}: invio di pochi bytes;
  \item dopo l' \textbf{ACK}, raddoppio la quantità di datagrammi inviati;
  \item quindi se i dati inviati diventano troppo grandi, si torna ad un
    passaggio indietro e si aumentano a passi i dati inviati.
\end{enumerate}

\section{Application layer}
L'ultimo livello di \emph{networking} è l'implementazione a livello
applicativo.
Per cui si rende necessaria un interfaccia per applicazione, \textbf{API}, si
userà lo standard \textbf{POSIX}; si utilizzano i \textbf{socket}.

\paragraph{File descriptor:}
un processo che permette di interagire con il \textit{file system}, permettono
di stabilire un collegamento tra i buffer a livello aplicativo e quelli di
sistema.

Il file descriptor (\textbf{fd}) creato, per esempio, alla chiamata di \textit{open},
sarà un \textbf{int} che rappresenta l'indice del file nell'array dei file
in utilizzo.

Nel caso di errore, \emph{fd} sarà un numero negativo.

\subsection{System calls sui socket di tipo datagram}
Ogni \textbf{system call} è bloccante, nel caso di un errore il programma si
ferma.
Si possono riprogrammare le \emph{syscall} per renderle non bloccanti, con la
gestione del messaggio di errore \textbf{E\_WOULDBLOCK}.

\subsubsection{Socket system call}
Quindi alla \textit{syscall} su \textbf{socket}, verrà creato un nuovo \emph{
fd}.

\begin{lstlisting}[caption=Esempio di syscall a socket]
int fd = socket(int domain, int type, int protocol);
\end{lstlisting}

Dove:
\begin{itemize}
  \item \textbf{domain:}
    \begin{itemize}
      \item \emph{AF\_UNIX}: locale, nomi di file;
      \item \emph{AF\_INET}: $IPV4$;
      \item \emph{AF\_INET6}: $IPV6$;
    \end{itemize}
  \item \textbf{type:}
    \begin{itemize}
      \item \emph{SOCK\_DGRAM}: $IPv4$;
      \item \emph{SOCK\_STREAM}: $IPv6$;
      \item \emph{SOCK\_RAW}: utilizzato per saltare i passaggi di trasporto e
        di rete fino a \emph{datalink};
    \end{itemize}
  \item \textbf{protocol:} esiste uno ed uno solo per type, è possibile
    lasciare il predefinito ($0$).
\end{itemize}

\subsubsection{Bind system call}
La \emph{syscall} \textbf{bind} è necessaria per associare un indirizzo ad un
\textbf{socket}.

\begin{lstlisting}[caption=Esempio di syscall a bind]
int fd = bind(int fd, const struct sockaddr*addr, socklen_t addrlen);
\end{lstlisting}

Dove:
\begin{itemize}
  \item \textbf{fd:} il file descriptor del \emph{socket} riferimento;
  \item \textbf{addr:} alloca lo spazio necessario per salvare il numero
    di $bit$ corretti;
  \item \textbf{addrlen:} la lunghezza dell'indirizzo da salvare (o del nome
    del file).
\end{itemize}

\subsubsection{Send to system call}
\begin{lstlisting}[caption=Esempio di syscall sendto]
ssize_t sendto(int fdA, const void*buf, size_t len, int flags, const struct sockaddr*addr, socklen_t addrlen);
\end{lstlisting}

Dove:
\begin{itemize}
  \item \textbf{fdA:} file descriptor del socket da cui inviare;
  \item \textbf{buf:} buffer da cui inviare;
  \item \textbf{len:} lunghezza del messaggio;
  \item \textbf{flags:} ;
  \item \textbf{sockaddr:} indirizzo del socket (comprensivo di porta);
  \item \textbf{socklen:} lunghezza del socket;
\end{itemize}

\subsubsection{Receive from system call}
\begin{lstlisting}[caption=Esempio di syscall receivefrom]
ssize_t recvfrom(int fdB, const void*buf, size_t len, int flags, struct sockaddr*addr, socklen_t addrlen);
\end{lstlisting}

Dove:
\begin{itemize}
  \item \textbf{fdA:} file descriptor del socket in cui ricevere;
  \item \textbf{buf:} buffer nel quale scrivere il messaggio;
  \item \textbf{len:} lunghezza del messaggio;
  \item \textbf{flags:} ;
  \item \textbf{sockaddr:} indirizzo del socket (comprensivo di porta del
    sender);
  \item \textbf{socklen:} lunghezza del socket;
\end{itemize}

\subsubsection{Connect system call}
\begin{lstlisting}[caption=Esempio di syscall connect]
int recvfrom(int fdA, const struct sockaddr*addr, socklen_t len);
\end{lstlisting}

\subsubsection{Listen system call}
\begin{lstlisting}[caption=Esempio di syscall listen]
int recvfrom(int fdB, int backlog);
\end{lstlisting}

\subsubsection{Accept system call}
\begin{lstlisting}[caption=Esempio di syscall accept]
int accept(int fdB, 2 parametri);
\end{lstlisting}
Il ritorno della funzione è un \textbf{fd}.
Per ogni successo il server associa un nuovo socket.

\subsection{Esempio di scambio di datagrammi}
Il server ed il client creano un socket di tipo \emph{datagram}, il server
esegue un \emph{bind}.
Il client invece esegue una \emph{syscall} di tipo \textbf{sendto}, il server
utilizza \textbf{recvfrom}, etc..

\subsection{Esempio di scambio di stream}
Il server ed il client creano socket di tipo \emph{stream}, il server esegue
la \emph{bind}, il client la \emph{connect}, per ogni \emph{send} e
\emph{receive} non è più necessario lo scambio di indirizzo.
Il server deve attivare \textbf{listen} prima dell'arrivo della \emph{connect}.

\section{Modello banda latenza}
$n=$numero di byte da inviare
\[
  D(n)=L_0+\dfrac{n}{B}
\]

\section{Protocollo a livello applicativo: DNS}
Lo scopo è di trasformare gli indirizzi numerici in stringhe.
Lo scopo della funzione \emph{addrinfo()} è di trasformare appunto l'indirizzo.

Il protocollo prevede che ci sia un \emph{client} che invia una richiesta ed
un \emph{server} che risponda, basato sul protocollo di trasporto \textbf{UDP},
è possibie utilizzare il protocollo \textbf{TCP} per poter inviare datagrammi
più estesi.

La porta standard è la $53$.

I datagrammi (serie da $32bit$, con dimensione massima di $512byte$) inviati
sono composti da:

\begin{tabular}{|c|c|}
                                               \hline
  ident (16bit)        & FLAG(16bit)        \\ \hline
  \#domande             & \#risposte          \\ \hline
  \#record autoritativi & \#record aggiuntivi \\ \hline
                       &                    \\ \hline
                       &                    \\ \hline
                       &                    \\ \hline
                       &                    \\ \hline
                       &                    \\ \hline
\end{tabular}

Il \emph{\#domande} indica quante righe precedono le \emph{risposte}.
I record autoratitativi seguono le risposte, a loro volta sono seguiti dai
dai record aggiuntivi.

\subsection{Tipi di indirizzi}
\begin{itemize}
  \item \textbf{Tipo A:} indirizzo standard (es. www.dibris.unige.it);
  \item \textbf{Tipo MX:} indirizzo del server di posta elettronica;
  \item \textbf{Tipo NS:} informazione interna del protocollo \emph{DNS},
    simile a \textbf{MX}, indica il server \textbf{DNS} del dominio;
  \item \textbf{Tipo CNAME:} la \textbf{C} rappresenta \textbf{Canonical},
    possono corrispondere più indirizzi alla stessa macchina, \emph{CNAME}
    serve per identificare la macchina interpellata.
\end{itemize}

\subsection{Risoluzione dei nomi}
Per ottenere l'indirizzo di un qualunque server collegato alla rete si
utilizza una struttura gerarchica.

\begin{enumerate}
  \item \textbf{Root}: conosce tutti i domini \emph{top level};
  \item \textbf{Top level:} i domini top level gerarchicamente sono i figli di
    root, esempio: \emph{it,com,uk etc.};
  \item \textbf{Autoritativi:} figli di \emph{top level}, esempio: \emph{unige,
    google, amazon}, un livello \emph{autoritativo} può avere come figlio
    un altro livello \emph{autoratitativo}, fino a circa $127$ livelli;
\end{enumerate}

\subsection{Algoritmo ricorsivo}
Il client invia una richiesta al suo server in locale, il quale va a vedere
l'indirizzo, quindi contatta il server root perchè non ha abbastanza
informazioni per risolvere la richiesta.

Quindi il server \emph{root} inoltra la stessa richiesta al \emph{top level}
domain corretto, il quale a propria volta invia la medesima richiesta al server
\emph{autoritativo} e così via.

Una volta terminata la ricorsione l'indirizzo del server contenuto nell'ultimo
dominio autoratitativo viene mandato indietro fino al client seguendo la
chiamata ricorsiva (risposta autoritativa).

\subsection{Algoritmo iterativo}
Il client invia la richiesta al suo server in locale, che se non riesce a
rispondere, il client reindirizza la richiesta al server root, questo rimbalzo
si ripete fintanto che il dominio corretto non risponde (risposta
autoritativa).

\subsection{Algoritmo ibrido}
I server che contengono le informazioni sono contattati in maniera iterativa,
i local servers invece supportano l'algoritmo ricorsivo.

Il server in locale piò dare una risposta non autoritativa se il nome è
salvato in cache.

\section{Protocollo NTP}
Il protocollo \textbf{NTP} (Network Time Protocol) si basa sulla gestione del
tempo, necessario per la sincronizzazione di orologi nella rete.

La versione attuale è la $4$ ed è definita nel protocollo $RFC5905$.


Il problema che affronta questo protocollo è relativo agli orologi locali dei
calcolatori che si connettono tra di loro.


Il protocollo \textbf{NTP} utilizza $4$ timestamp:
\begin{itemize}
  \item $t_1$: il tempo del momento in cui viene inviato il messaggio
    (client);
  \item $t_2$: il tempo del momento in cui viene ricevuto il messaggio
    (host);
  \item $t_3$: il tempo del momento in cui viene inviato il messaggio
    (host);
  \item $t_4$: il tempo del momento in cui viene ricevuto il messaggio
    (client).
\end{itemize}

Quindi a differenza dei protocolli \textbf{UTP} e \textbf{TCP} individuo $2$
timestamp in più.

Il \emph{Round Trip Time} sarà dato da:
\[
  \text{RTT}=(t_4-t_1)-(t_3-t_2)
\]


La \textbf{epoca} è lo standard cui si fa riferimento secondo cui dal punto
di vista dei calcolatori comincia il tempo (1 gennaio 1970).

Considerando la precisione teorica degli orologi a disposizione dei calcolatori
, i server \textbf{NTP} hanno una struttura gerarchica:
\begin{itemize}
  \item \textbf{livello 0:} orologio atomico;
  \item \textbf{livello 1:} macchina collegata direttamente allo strato $0$;
  \item \textbf{livello 2:} macchine collegato allo strato $1$, i livelli
    successivi sono ricorsivi.
\end{itemize}
Ognuna di queste macchine è sia client, che server.

\chapter{Posta elettronica}
Una delle caratteristiche principali della posta elettronica è l' \textit{
asincronicità}.

\section{SMTP}
Le email utilizzano diverse protocolli tra cui \textbf{SMTP}, definito per la
prima volta nell' \textbf{RFC 821} (che risale al 1980) ed espanso nell'
\textbf{RFC 5321}.

Utilizza una codifica \textbf{ASCII a 7bit}, la porta predefinita è la $25$ ed
utilizza il protocollo \textbf{TCP}.

\subsection{Protocollo Best Effort}
L'applicazione del protocollo \textit{Best Effort} nell'ambito del protocollo
\emph{SMTP} implica un salvataggio da parte del mittente dei dati da inviare
fintanto che non si è certi che il messaggio sia arrivato, oppure alla
scadenza di un timeout.

\subsection{Funzionamento}
Il mittente (client) invia il messaggio ad un server, quindi il client alla
fine dell'invio può chiudere la connessione.
A questo punto il server che ha ricevuto il messaggio lo invia ad un altro
server, se necessario, con l'obiettivo di spostare il messaggio al server
del ricevente, in questo punto termina il compito del protocollo \textbf{SMTP},
successivamente il ricevente riceve il messaggio nel momento
in cui apre una connessione con il server.

\subsubsection{Three way connection}
Usa la \emph{three way connection}, una volta stabilita la connessione (dopo il
segnale di ACK)

\begin{enumerate}
  \item \textbf{SERVER:} invia la stringa \emph{220 server@example.com};
  \item \textbf{CLIENT:} invia la stringa \emph{HELO server@example.it}, oppure
    la stringa \emph{EHLO server@example.it} per connessioni sicure;
  \item \textbf{SERVER:} invia la stringa \emph{250 client@example.it};
  \item \textbf{CLIENT:} invia la stringa \emph{MAIL FROM: client@example.it};
  \item \textbf{SERVER:} invia la stringa \emph{250 client@example.it server
    ok};
  \item \textbf{CLIENT:} invia la stringa \emph{RECPT TO: recipient@example.it}
    ;
  \item \textbf{SERVER:} invia la stringa \emph{250 recipient@example.it 
    recipient OK};
  \item \textbf{CLIENT:} invia la stringa contenente \emph{DATA};
  \item \textbf{SERVER:} risponde con $354$;
  \item \textbf{CLIENT:} invia il testo del \textbf{messaggio} che si conclude
    con un newline ed un punto;
  \item \textbf{SERVER:} risponde con \emph{250} segnalando la fine della
    comunicazione;
  \item \textbf{CLIENT:} invia la stringa \emph{quit}.
\end{enumerate}

\subsubsection{Formato del messaggio}
Il messaggio è composto da $72$ caratteri per riga, le sequenze di caratteri
\emph{<CR><LF>} costituiscono un a capo; lo \emph{header} termina con due
a capi di seguito.

Lo \emph{header} è composto di chiavi, alcune sono ridondanti rispetto ai dati
scambiati secondo il protocollo \textbf{SMTP}, difatti i dati esposti nel
corpo del messaggio sono puramente a scopo informativo.
\begin{itemize}
  \item subject;
  \item from;
  \item to;
  \item cc;
  \item bcc;
  \item date
\end{itemize}


I messaggi possono essere codificati in base $64$ Radix, come caratteri
comprende le cifre, le lettere maiuscole e minuscole, $\_$ ed $=$.

Quindi ad ogni $6bit$ del messaggio corrisponde $1$ carattere Radix.

La prima versione del protocollo \textbf{SMTP} utilizzava la codifica \emph{
uuencode}.


Un evoluzione del Radix 64 utilizza i MIME, per esempio un \emph{MIME/TYPE}
\emph{text/ascii} indica che devo codificare un file di testo formato ascii.

\textbf{MAIL USER AGENT}

\textbf{POP3}

\textbf{TCP/IP}

\textbf{IMAP}
